% Proposed experimental protocol, not a report of completed experiments.
% Requires amsmath. Replace proposed defaults only after development validation.
% Source links and all engineering details are in SOURCES.txt and AGENT_INSTRUCTIONS.txt.

\subsection{Downstream Evaluation}
\label{sec:downstream_setting}

\paragraph{Scope.}
We propose two complementary studies: sustained execution of frozen robot
policies in generated environments, and control-conditioned video prediction
with target-embodiment-adapted world models. The first study tests manipulation
and mobile manipulation over persistent task streams. The second measures
prediction quality, task-stage outcome accuracy, and persistence through
occlusion. All settings below are preregistration defaults; the resulting
scores must be computed from completed, audited runs.

\paragraph{World construction and splits.}
For each robot setting, we use 30 held-out world bundles and two rollout seeds
per bundle. Each lifetime contains ten scheduled instructions. Manipulation
uses a fixed-base dual-ARX-X5 workspace. Mobile manipulation uses a Lift2
platform with R5a arms and a connected ten-room world: task one uses the seed
room, and each subsequent task requires an appended room. The room sequence,
object instances, task goals, and initial snapshots are shared across policies.
The two seeds retain the same generated geometry.

The worlds are stratified across delivery/recollection,
sorting/reorganization, and multi-object staging/cleanup, with ten bundles
per family. Eight separate development worlds per setting are reserved for
interface verification and protocol calibration. World groups and exact
generated asset instances are disjoint between training, development, and
test; derived snapshots and action branches remain in their source group.
This split is relative to the present study, without assuming knowledge of
foundation-model pretraining corpora.

Generated assets are compiled into the native simulators with persistent
instance identities, validated collision geometry, and declared physical
parameters. Admission verifies stable placement, robot clearance,
reachability, traversable portals, executable goals, and a physically
executed reference sequence. Admission does not depend on the tested
policies. We report generation yield separately. Primary tasks use rigid
objects and open receptacles; articulation and deformable interaction require
separately validated extensions.

\paragraph{Persistent robot deployment.}
We evaluate the native compatible checkpoints of $\pi_{0.5}$, GR00T N1.7,
InternVLA-A1, OpenWAM-$\alpha$, and InternW0-$\Delta$ for manipulation, and
$\pi_{0.5}$, $\pi_0$, InternVLA-A1, and OpenWAM-$\alpha$ for mobile
manipulation. Exact checkpoint, source, observation, action, and controller
manifests are frozen before testing. No robot-policy training on WorldEdge
is included. A row is reportable only after its actual interface is validated.

A lifetime is reset only at its beginning. Between instructions, all actual
robot and object states persist, including unsuccessful placements, contact
states, and velocities. Room insertion preserves all existing entities.
Queued actions for the old instruction are discarded, while permitted native
policy memory is retained. Policies observe the current instruction and
their declared native sensor inputs, without privileged scene state or
future instructions. Task failure does not terminate an otherwise usable
lifetime. Unreachable or unattempted goals remain in the scheduled
denominator, and no object is respawned or replaced.

We allocate 120 simulated seconds per manipulation task and 300 seconds per
mobile task. These deadlines and physical predicate tolerances are calibrated
only on development worlds. The offline capability test pauses physical time
during model inference; inference latency and total wall time are reported
separately. Native action timing is recorded per embodiment and is not
conflated with action-chunk count. Success is computed independently from
simulator state and event traces, requiring the specified final conjunction,
temporal ordering, object release, and a two-second stable dwell.

\paragraph{Robot metrics.}
Let $y_{i,t}\in\{0,1\}$ denote success for task $t$ in lifetime $i$, with
unreached tasks assigned zero. Strict chain success and completed tasks are
\begin{align}
\mathrm{SR@}K &=
\frac{100}{N}\sum_{i=1}^{N}\prod_{t=1}^{K}y_{i,t},
\qquad K\in\{5,10\},\\
\mathrm{Tasks}/10 &=
\frac{1}{N}\sum_{i=1}^{N}\sum_{t=1}^{10}y_{i,t}.
\end{align}
The second metric includes successful recovery after an earlier failure.
For identity-specific anchor/reuse pairs
$\mathcal{P}=\{(1,4),(4,7),(7,10)\}$, we report
\begin{equation}
\mathrm{Reuse}=
\frac{100}{3N}\sum_{i=1}^{N}
\sum_{(a,b)\in\mathcal{P}}y_{i,a}y_{i,b}.
\end{equation}
Anchor and reuse tasks require actual displacement and correct-instance
placement. The denominator includes every scheduled pair. Stationary reuse
measures sustained instance-specific interaction; mobile reuse additionally
requires returning to or searching previously encountered spaces.

\paragraph{Video world models and data.}
The primary video comparison uses Cosmos-Predict2.5-2B with its
robot/action-conditioned interface and DreamDojo-2B. DreamDojo-14B is an
optional capacity extension. All primary methods are adapted to the same
dual-ARX embodiment using the same training split, then frozen.
Unrelated native robot checkpoints are not pooled into this comparison.
Mobile video prediction requires a separate adaptation and reporting panel.

We construct 3,200 training clips from 160 worlds, 400 development clips from
20 worlds, and 600 test clips from the same 30 manipulation test worlds used
for policy evaluation. Each world contributes ten matched success/failure
pairs: three placement, two sorting, two multi-object ordering, and three
revisit pairs. Paired branches share an initial snapshot, camera, stage goal,
and physical parameters. Outcomes are established by actual simulator
execution independently of the video models. The test set therefore contains
300 positive and 300 negative stage outcomes and 180 revisit clips.

Each clip lasts 30 seconds at 10 Hz. A dedicated collector executes
piecewise-constant full bimanual commands for 100 ms each, using 25 physics
substeps at 250 Hz. Any resampled policy controls are re-executed to generate
their corresponding reference; old videos are not reused with modified
actions. Model inputs are the initial RGB image and issued controls.
Both predictors use the same fixed empty text prompt.
Future measured robot poses, future object states, and latent actions
extracted from reference video are excluded. No room is inserted during a
scored clip.

\paragraph{Adaptation and prediction.}
We propose 20,000 target-adaptation updates per model, an effective batch of
32, 13-frame samples, $480\times640$ images, BF16, and AdamW with initial
learning rate $10^{-5}$, 1,000 warm-up updates, cosine decay to $10^{-6}$,
weight decay 0.01, and gradient clipping at 1.0. Video-transformer and
action-conditioning parameters are trained; the text encoder and video codec
are frozen. Each method's native loss is recorded. The protocol matches data,
updates, and optimizer schedule rather than computational cost or objective
functions. Training uses seed 42; test uncertainty is conditional on the
selected adapted checkpoint.

Checkpoints at 5k, 10k, and 20k updates are selected by mean endpoint LPIPS
at 30 seconds on a fixed 100-clip development subset. Test data do not enter
training or checkpoint selection. Each test clip is predicted with seeds
zero and one, without best-of-sample selection. A 13-frame recursive interface
adds twelve new frames per chunk; 25 chunks produce 301 timestamped frames
including the initial observation. Continuation uses generated context only,
with no intermediate reference-image refresh.

\paragraph{Video metrics.}
LPIPS uses the fixed AlexNet implementation and version 0.1, with identical
RGB normalization and resizing to $240\times320$. We compare frames at
exactly 5 and 30 seconds, without temporal warping. These are endpoint
distances. An independent blinded visual reader classifies the outcome of a
predeclared 30-second stage goal. Success requires the terminal predicates
during the last 0.5 seconds and any required event order; it is not the
eventual outcome of an arbitrary longer parent task. Outcome is balanced
accuracy:
\begin{equation}
\mathrm{Outcome}=
50\left(
\frac{\mathrm{TP}}{N_+}+
\frac{\mathrm{TN}}{N_-}
\right),
\end{equation}
where generation samples are averaged within each clip and unreadable
predictions are incorrect for both classes.

Revisit cases establish target identity, physically change its state during
approximately 2--8 seconds, occlude it for at least five seconds, and reveal it
during 25--30 seconds. Both branches change the actual target state.
The three revisit pairs per world target 5, 10, and 15 seconds of occlusion.
The reveal suffix is identical across paired branches and does not reach
toward the hidden target. Revisit requires the correct identity and actual
visible state at the endpoint, including the actual wrong placement in a
failed task branch. Missing, duplicated, or unreadable targets are incorrect.

The visual reader is calibrated on 200 balanced development clips against
simulator labels, with human auditing. We require at least 95\% balanced
accuracy and 90\% sensitivity/specificity before automatic test scoring.
All unknown outputs and a stratified 20\% of generated videos are audited;
categories with over 10\% disagreement receive complete adjudication.
Static, looping, stale-state, and corrupted-action controls are reported in
the supplement. Missing infrastructure outputs remain scheduled failures for
semantic scores, and an incomplete LPIPS horizon is marked not reportable.

\paragraph{Statistics and interpretation.}
We bootstrap complete world groups 10,000 times, retaining all tasks,
paired branches, and repeated seeds within each group. The same resamples
are used for paired method comparisons; policy resampling is stratified by
task family. Main tables show point estimates, while the supplement provides
95\% intervals, event counts, coverage, runtime, and failure causes.
The full nine-row policy schedule contains 540 lifetimes and at most 5,400
task attempts. The two primary video models produce 2,400 test predictions.

The policy table uses a frozen task/growth schedule for comparability.
A separate optional demonstration may generate tasks adaptively from actual
success/failure feedback. Its success rates are not mixed with the matched
table. Control-conditioned video replay measures prediction of an issued
trajectory; it does not by itself establish planning benefits or policy
execution in an imagined environment.
